\documentclass[11pt,a4paper]{article}
\usepackage[margin=2.4cm]{geometry}
\usepackage{graphicx,booktabs,microtype,amsmath}
\usepackage[svgnames]{xcolor}
\usepackage[colorlinks=true,linkcolor=MidnightBlue,urlcolor=MidnightBlue]{hyperref}
\usepackage{parskip}
\emergencystretch=2.2em \tolerance=1200
\usepackage{titlesec}
\titleformat{\section}{\large\bfseries}{\thesection}{0.7em}{}
\titleformat{\subsection}{\normalsize\bfseries}{\thesubsection}{0.6em}{}
\newcommand{\hco}{HCO$_3^-$}
\newcommand{\cotwo}{CO$_2$}
\usepackage{tcolorbox}

\title{\vspace{-1.5cm}A drainage-concentration ceiling on enhanced rock weathering\\
\large Request for comment on a bound we have implemented but not shipped}
\author{Zeke Hausfather\\\small Stripe / Frontier}
\date{\today}

\begin{document}
\maketitle
\vspace{-0.7cm}

\begin{tcolorbox}[colback=Gainsboro!25,colframe=Gray,boxrule=0.4pt,arc=2pt]
\textbf{What we are asking.} We have built a global gridded ERW suitability model.
While auditing it we found it reports carbon that could not physically leave the
field dissolved in the field's own drainage water. We implemented a bound, and it
changes the map a great deal --- including one of our headline results. Before we
ship it we would like the ERW community to tell us whether the bound is right, and
in particular whether the questions in \S\ref{sec:questions} have answers we
have missed. The bound is currently switched \emph{off} in the published map.
\end{tcolorbox}

\section{The problem}

Our model, like most process-based ERW models, computes a mineral dissolution rate
and converts dissolved rock to \cotwo{} removal. Nothing in that chain requires the
resulting carbon to be carryable at a chemically possible bicarbonate concentration.

Dividing our own \cotwo{} flux by our own water flux, area-weighted over global
cropland at 30\,t\,ha$^{-1}$\,yr$^{-1}$ of basalt, the
model implies a drainage \hco{} concentration of
\textbf{8\,/\,20\,/\,74\,mmol\,L$^{-1}$} at p10\,/\,p50\,/\,p90.
For scale, world rivers run 0.5--4\,mmol\,L$^{-1}$.

This is not specific to our model. Beerling et al.\ (2024) report a CDR$_{\rm pot}$
of 10.5\,t\,\cotwo{}\,ha$^{-1}$ over four years at the Illinois Energy Farm and
report no drainage chemistry; at Illinois tile drainage of 200\,mm\,yr$^{-1}$ that
requires 29.8\,mmol\,L$^{-1}$. Kelland et al.\ (2020) closes the loop inside one
experiment: measured cation release over measured drainage requires
21.1\,mmol\,L$^{-1}$, while measured leachate alkalinity was
1.10\,$\pm$\,0.147\,mmol\,L$^{-1}$ and statistically indistinguishable from control
--- a 19$\times$ shortfall measured on both sides.

\section{Which water flux is $q$?}
\label{sec:drainage}

The bound is $q \cdot [\mathrm{HCO_3^-}]_{\max} \cdot 44$, so it is linear in the
water flux. The choice of flux therefore sets the level as directly as the chemistry
does, and we would like it checked alongside the chemistry.

WaterGAP2-2e publishes four candidates, and over global cropland they differ by a
factor of five (area-weighted median, mm\,yr$^{-1}$): groundwater recharge $q_r$
75, subsurface runoff $q_{sb}$ 75, surface runoff $q_s$
84, total runoff $q_{tot}$ 178.

The intuitive choice is $q_r$: ERW needs water percolating below the root zone, not
overland flow. It cannot be used on its own. $q_r$ is \emph{exactly zero} over 0.10\%
of cropland area, concentrated in river deltas --- 23\% of the Mekong Delta's
cropland, 24\% of the Red River delta --- where the water table is at the surface.
WaterGAP is right that nothing recharges an aquifer there, but field drainage still
leaves laterally to canals with its bicarbonate in it, so those cells come out at
zero ERW potential in some of the wettest cropland on Earth.

$q_{sb}$ is not the fix, though it looks like it. In WaterGAP recharge feeds the
groundwater store and that store discharges as baseflow, so a 30-year mean $q_{sb}$
is close to $q_r$ relabelled (global land medians 33.5 vs 32.6\,mm\,yr$^{-1}$). It
clears the deltas and creates worse zeros where groundwater is pumped: 26\% of the
Indo-Gangetic Plain.

We use $q_{tot}$, on the argument that Maher \& Chamberlain fit $D_w$ against
catchment discharge per unit area --- which \emph{is} total runoff --- so driving a
$q_{tot}$-calibrated $D_w$ with recharge penalises the flux twice. The
counter-argument we cannot dismiss is that surface runoff has little contact time
with topsoil rock, so $q_{tot}$ credits water that may have weathered nothing; our
reply is that an effective catchment-scale $D_w$ already absorbs that. Treat the two
as a bracket: unbounded global gross 2.49 on $q_{tot}$ against 2.15 on $q_r$,
and bounded 0.91 against 0.36. \textbf{This is question 5 in
\S\ref{sec:questions}.}

\section{A wrong first answer, because it is instructive}

Our first attempt computed the ceiling as $q\,[\mathrm{HCO_3^-}](\mathrm{pH},
p\mathrm{CO_2})$ holding each cell's \emph{pre-treatment} soil pH fixed. That gives
0.42\,mmol\,L$^{-1}$ at the median and an apparent 563$\times$ exceedance.

It is wrong, and wrong in an interesting way. pH is not exogenous: adding base
cations at fixed $p$\cotwo{} raises alkalinity and pH \emph{together}, and raising pH
is precisely what a silicate amendment does. Reading the carbonate equation with pH
held fixed answers a different question. Tellingly, 0.42\,mmol\,L$^{-1}$ is, to two
significant figures, the observed mean alkalinity of streams draining
\emph{unamended} volcanic rock --- a good baseline, and the wrong ceiling.

\section{The bound we propose}

The bound is where the rising pH meets \textbf{carbonate saturation}. Solving charge
balance $2[\mathrm{Ca}] + 2[\mathrm{Mg}] = [\mathrm{HCO_3^-}]$ simultaneously with
fixed $p$\cotwo{} and a calcite saturation state $\Omega$ gives a closed form:

\begin{equation}
[\mathrm{HCO_3^-}]_{\max} = \left(
\frac{2\,\Omega\,K_1 K_H\, p\mathrm{CO_2}\, K_{\rm sp}}{f_{\rm Ca} K_2}
\right)^{1/3}
\end{equation}

with $f_{\rm Ca}$ the Ca share of the divalent-cation charge (only Ca constrains
calcite; magnesite is kinetically inhibited at surface temperature, so a lower
$f_{\rm Ca}$ \emph{raises} the ceiling). Gross CDR is then capped at
$q\,[\mathrm{HCO_3^-}]_{\max}\,M_{\rm CO_2}$. $K_1$, $K_2$, $K_H$ and $K_{\rm sp}$
are Plummer \& Busenberg (1982), evaluated at each cell's own temperature.

The cube root makes this robust: being wrong about soil $p$\cotwo{} by 5$\times$ moves
the ceiling only 1.7$\times$. At the protocol-mandated
4,000\,$\mu$atm, 15\,\textdegree C and $f_{\rm Ca} =
0.5$, it gives \textbf{3.03\,mmol\,L$^{-1}$ at $\Omega = 1$}
and \textbf{6.52\,at $\Omega = 10$}.

We adopt $\Omega = 10$ because carbonate
precipitation is kinetically inhibited by DOC and phosphate. Zhang et al.\ (2022)
state precipitation in river water ``is generally observed to be negligible at
$\Omega < 10$'' and run their own model over $\Omega = 5$--25; soils carry far more
DOC than rivers. \textbf{This is the single largest discretionary choice in the
bound} --- see question 2.

\subsection*{Applied to the carbon, not to the rock}

We cap the \cotwo{} and leave the predicted \emph{fraction weathered} unbounded.
Rock can dissolve without the carbon leaving, and the gap between those two layers
is then a visible, meaningful quantity rather than an inconsistency. The cap is
applied \emph{after} the alkalinity-to-DIC conversion, because what has to fit in
the water is the bicarbonate.

\section{Validation}

\textbf{Against a textbook benchmark.} Setting $f_{\rm Ca} = 1$, $\Omega = 1$ and
$p$\cotwo{} to atmospheric should reproduce the classic open-system calcite
equilibrium. It gives 0.949\,mmol\,L$^{-1}$ alkalinity,
0.475\,mmol\,L$^{-1}$ Ca and pH 8.20, against the textbook
$\sim$1.0, $\sim$0.5 and $\sim$8.3. Neglecting activity coefficients and the
CaHCO$_3^+$ ion pair biases this low by 10--20\%, so our ceiling is mildly
conservative.

\textbf{Against five independent measurements} that share no assumptions with the
closed form (Figure~\ref{fig:main}A):

\begin{center}
\begin{tabular}{lr}
\toprule
Anchor & mmol\,L$^{-1}$ \\
\midrule
Riverine transport capacity, Zhang et al.\ 2022, back-converted & 4.3--13 \\
Midwest agricultural tile drainage and porewater, Hamilton et al.\ 2007 & 1--7 \\
World pristine rivers, 99th percentile (Meybeck) & 5.95 \\
Carbonate-terrain streams (Meybeck) & 3.15 \\
\midrule
\textit{Achieved} in ERW field trials --- an observation, not a ceiling & 0.11--0.75 \\
\bottomrule
\end{tabular}
\end{center}

The closed form lands inside all four bounds, and straddles the most relevant one
(agricultural tile drainage). Note the last row: measured ERW trials \emph{achieve}
9--60$\times$ \emph{below} the ceiling. The bound is a rail
that makes an impossible claim impossible; it is not why modelled ERW rates exceed
measured ones, which is a cation-retention and lab-to-field-rate problem.

\begin{figure}[t]
\centering
\includegraphics[width=\linewidth]{rfc_flux_reconciliation_fig.pdf}
\caption{\textbf{A} What the model requires against what the water can hold, with
five independent anchors on one axis. \textbf{B} The consequence that matters: the
ceiling falls with warming while the rate law rises, so imposing it cuts the
warm-climate gradient from 3.95$\times$ to 1.41$\times$ rather than merely
rescaling the level. Area-weighted medians over global cropland at
30\,t\,ha$^{-1}$\,yr$^{-1}$.}
\label{fig:main}
\end{figure}

\section{Consequences}

The ceiling binds on \textbf{93.0\% of cropland area}. Median gross CDR
falls from 1.589 to 0.510\,t\,\cotwo{}\,ha$^{-1}$\,yr$^{-1}$; the global
total from 2.488 to 0.910\,Gt\,\cotwo{}\,yr$^{-1}$. The $\Omega$ choice
spans 0.240 ($\Omega = 1$) to 0.510 ($\Omega = 10$) at
the median, and that spread is the honest uncertainty on the level.

\textbf{But the level is not the interesting part.} $C_{\rm eq}$ \emph{falls} with
warming (3.56\,/\,3.03\,/\,2.58\,mmol\,L$^{-1}$ at 5\,/\,15\,/\,25\,\textdegree C)
while the rate law rises steeply, so the exceedance is monotonic in temperature:

\begin{center}
\begin{tabular}{lrrrr}
\toprule
Mean soil T (\textdegree C) & \% of area & Unbounded & Ceiling & Exceedance \\
\midrule
0--10 & 28.3 & 0.777 & 0.410 & 1.9$\times$ \\
10--15 & 19.0 & 1.058 & 0.600 & 1.8$\times$ \\
15--20 & 15.2 & 1.566 & 0.447 & 3.5$\times$ \\
20--25 & 15.5 & 3.037 & 0.711 & 4.3$\times$ \\
25--45 & 21.9 & 3.067 & 0.578 & 5.3$\times$ \\
\bottomrule
\end{tabular}
\end{center}

\noindent The warmest-to-coolest ratio of the median goes from
\textbf{3.95$\times$ unbounded to 1.41$\times$ at the ceiling}. Imposing the
bound therefore largely \emph{removes} the warm-climate advantage: most of the 3.95$\times$ advantage an unbounded rate law produces is not supported by the water. That is the
most consequential and least intuitive result here, and the one we would most like
checked.

\textbf{It also decouples carbon from application rate.} The ceiling depends on
drainage and carbonate chemistry, not on how much rock is on the field. Raising our
nominal rate from 20 to 30\,t\,ha$^{-1}$ raised the
unbounded median 50\% and the bounded global total by 8.6\%. Past the point where
drainage saturates, extra feedstock raises the transport-limited share of global
cropland rather than the tonnage.

\section{Two points from Maher \& Chamberlain, for the record}

Working the primary source resolved two things that may be of wider interest.

\textbf{$\tau = e^2$ is real and is not already in the Fig.~2 contours.} Their Eq.~3
is $C = C_{\rm eq}(\tau D_w/q)/(1 + \tau D_w/q)$, with $\tau$ in Table~S1. The Fig.~2
contour \emph{labels} reproduce to two significant figures from the paper's own
printed $R_{n,\max}$, $C_{\rm eq}$ and $L_\phi$ using bare $D_w = L_\phi/T_{\rm eq}$,
while Fig.~2B's plotted flux plateaus match $C_{\rm eq}\tau D_w$ and are 0.8 of a
decade off without $\tau$.

\textbf{Consequently all cropland is on the low-$q$ limb.} $\tau D_w$ is
222\,mm\,yr$^{-1}$ at $D_w = 0.03$ and
2,217 at $D_w = 0.3$, both at or above the p90 of cropland drainage
(613\,mm\,yr$^{-1}$). On that limb the flux is $C_{\rm eq}q$ and kinetics drop
out --- Maher (2010, p.~104): beyond $L_{\rm eq}$ the flux ``conveys no information
on the actual weathering kinetics or available surface area.'' Godsey et al.\ (2009)
measure near-chemostatic concentration--discharge slopes of $-0.05$ to $-0.15$
across 59 catchments, which is the same statement empirically.

\section{Status}

The bound is implemented and tested, and we are \emph{not} currently applying it to
the \cotwo{} figures we report, pending this review. We do report the exceedance
alongside them: 93.0\% of cropland area exceeds the bound by a median
2.9$\times$, so the unbounded figures should be read as an upper limit on
dissolution rather than as carbon shown to leave the field.

On external consistency: unbounded, our global total is
2.49\,Gt\,\cotwo{}\,yr$^{-1}$; bounded, 0.91. Both sit inside the 0.5--4
range of published global ERW estimates, and we do not read that as support. Those
published estimates are not transport-bounded either, so agreement with them is not
evidence about this term.

\section{Questions we would like answered}
\label{sec:questions}

\begin{enumerate}
\item \textbf{Is carbonate saturation the right ceiling for \emph{alkalinity}?} Silica
is bounded by clay equilibrium, but alkalinity is not --- observed
\hco{}:SiO$_2$ ratios of 2.1--9.3 say the two are decoupled by incongruent
weathering. Everything else rests on this step, and we would most like it attacked.
\item \textbf{Is $\Omega = 10$ defensible in soil, or should it be $\Omega = 1$?} That
pair spans 0.240--0.510\,t\,\cotwo{}\,ha$^{-1}$\,yr$^{-1}$ at the median and we
know of no measurement that discriminates in an amended agricultural soil.
\item \textbf{Is capping the carbon while leaving fraction weathered unbounded the
right separation?}
\item \textbf{Does the temperature dependence survive scrutiny?} It is what collapses
the warm/cool ratio from 3.95$\times$ to 1.41$\times$, and it is the result
we are least confident about.
\item \textbf{Which water flux should drive the bound (\S\ref{sec:drainage})?}
Recharge and total runoff differ 2.4$\times$ over cropland and the bound is linear in
the flux. Recharge gives zero drainage across major river deltas, so it cannot be
right on its own; total runoff credits surface flow that had little contact with the
applied rock. If there is field evidence on what fraction of agricultural runoff
carries weathering products --- tile-drain versus overland partitioning of alkalinity
export --- that would settle the largest single uncertainty in this note.
\end{enumerate}

A sixth, if anyone has data: the bound is unvalidated on flooded/paddy soils, where
the mandated 50{,}000\,$\mu$atm lifts it to 13--18\,mmol\,L$^{-1}$, above every
anchor above. We could find no measured floodwater alkalinity or paddy lateral DIC
export flux anywhere in the literature.

%%% APPENDIX-A -- aggregated commercial deliveries. Remove this block if the
%%% sharing terms do not permit it; nothing else references it.
\appendix
\section{Aggregate check against verified commercial deliveries}

Applying the same test to eight independently verified commercial ERW deliveries
(basalt; roughly 15--100\,t\,ha$^{-1}$; India, US Corn Belt and Brazil), \textbf{seven
of the eight} report a CDR per hectare that would require more bicarbonate than
their own region's drainage can carry --- \textbf{0.6--7.8$\times$} the ceiling across
the set, \textbf{1.3--3.9$\times$} restricting to the three where CDR was measured
independently rather than derived from fraction weathered. The eighth is the wettest
site in the set at over 1{,}500\,mm\,yr$^{-1}$ of drainage, and its reported CDR is
carryable.

This is \emph{not} an over-crediting finding, and we would ask that it not be read
as one. Those figures are dissolution-based, so the comparison is ``how much rock
dissolved'' against ``how much carbon the water could carry out'', and both can hold
at once. The gap is then cation retention in secondary phases and export lag: 10--50$
\times$ more retained than exported across four soils and thirteen feedstocks
(Hammes et al.\ 2025), retarded fractions of 93--98\% (te Pas et al.\ 2025), and
modelled export lags of 5--22 years (Kanzaki et al.\ 2025) --- the same order as the
gap here. What it does establish is that dissolution-based CDR per hectare cannot be
read as export without a retention term.

Per-deployment values are withheld: they derive from independent verification
reports that are not ours to publish.
%%% END APPENDIX-A

\begin{thebibliography}{9}\small
\bibitem{beerling24} Beerling, D.J.\ et al.\ (2024) \textit{PNAS} 121, e2319436121.
\bibitem{godsey09} Godsey, S.E., Kirchner, J.W.\ \& Clow, D.W.\ (2009) \textit{Hydrol.\ Process.} 23, 1844--1864.
\bibitem{hamilton07} Hamilton, S.K.\ et al.\ (2007) \textit{Global Biogeochem.\ Cycles} 21, GB2021.
\bibitem{hammes25} Hammes, J.S.\ et al.\ (2025) EGUsphere preprint 2025-5402.
\bibitem{kanzaki25} Kanzaki, Y.\ et al.\ (2025) \textit{Environ.\ Res.\ Lett.} 20, 074055.
\bibitem{kelland20} Kelland, M.E.\ et al.\ (2020) \textit{Glob.\ Change Biol.} 26, 3658--3676.
\bibitem{maher10} Maher, K.\ (2010) \textit{Earth Planet.\ Sci.\ Lett.} 294, 101--110.
\bibitem{maher14} Maher, K.\ \& Chamberlain, C.P.\ (2014) \textit{Science} 343, 1502--1504.
\bibitem{pb82} Plummer, L.N.\ \& Busenberg, E.\ (1982) \textit{Geochim.\ Cosmochim.\ Acta} 46, 1011--1040.
\bibitem{tepas25} te Pas, E.E.E.M.\ et al.\ (2025) \textit{Front.\ Clim.} 6, 1524998.
\bibitem{zhang22} Zhang, S.\ et al.\ (2022) \textit{Limnol.\ Oceanogr.} 67, doi:10.1002/lno.12244.
\end{thebibliography}

\end{document}
